\documentclass[12pt]{article}
\usepackage[T1]{fontenc}
\usepackage[utf8]{inputenc}
\usepackage{lmodern}
\usepackage{textcomp}
\usepackage{enumitem}
\usepackage{geometry}
\geometry{margin=1in}
\begin{document}
\begin{center}
Answer Key - Psychology - Undergraduate Level Assessment 15 \
\small (Answers are ordered exactly as in the question paper.)
\end{center}

\section*{Multiple Choice}
\begin{enumerate}[label=\arabic*.]
\item \textbf{Answer:} A
\\ \textit{Options:} A) Perls.; B) Maslow.; C) Erikson.; D) Piaget.; E) Rogers.
\\ \textit{Note:} Correct option: A - Perls.
\item \textbf{Answer:} C
\\ \textit{Options:} A) too little curvature of the cornea and lens; B) too much curvature of the retina and lens; C) too much curvature of the cornea and lens; D) too much curvature of the iris and cornea; E) incorrect alignment of the cornea and retina
\\ \textit{Note:} Correct option: C - too much curvature of the cornea and lens
\item \textbf{Answer:} A
\\ \textit{Options:} A) implementation of the intervention by the researcher; B) pre-determined assignment of participants to groups; C) use of a non-comparison group; D) use of multiple measures of outcome; E) use of a single group of participants
\\ \textit{Note:} Correct option: A - implementation of the intervention by the researcher
\item \textbf{Answer:} A
\\ \textit{Options:} A) Triarchic theory of intelligence; B) Multiple intelligences theory; C) Cognitive development theory; D) The positive manifold; E) G theory
\\ \textit{Note:} Correct option: A - Triarchic theory of intelligence
\item \textbf{Answer:} B
\\ \textit{Options:} A) cell body; B) myelin; C) synaptic vesicles; D) glial cells; E) terminal buttons
\\ \textit{Note:} Correct option: B - myelin
\end{enumerate}

\section*{True / False}
\begin{enumerate}[label=\arabic*.]
\setcounter{enumi}{5}
\item \textbf{Answer:} True
\\ \textit{Options:} A) True; B) False
\\ \textit{Note:} LLM-generated true statement based on MCQ.
\item \textbf{Answer:} True
\\ \textit{Options:} A) True; B) False
\\ \textit{Note:} LLM-generated true statement based on MCQ.
\item \textbf{Answer:} True
\\ \textit{Options:} A) True; B) False
\\ \textit{Note:} LLM-generated true statement based on MCQ.
\item \textbf{Answer:} False
\\ \textit{Options:} A) True; B) False
\\ \textit{Note:} LLM-generated false statement. Correct answer: 'autonomy vs. shame and guilt'.
\item \textbf{Answer:} False
\\ \textit{Options:} A) True; B) False
\\ \textit{Note:} LLM-generated false statement. Correct answer: 'cognitive-dissonance'.
\end{enumerate}

\section*{Long Form Response}
\begin{enumerate}[label=\arabic*.]
\setcounter{enumi}{10}
\item \textbf{Answer:} An organism's behavior in relation to a goal with conflicting motives would vary depending on the distance from the goal and the relative strengths of the approach and avoidance drives.. Explain why this is the correct answer and provide supporting evidence. Reference key facts or concepts that support this choice.
\\ \textit{Note:} Long-form derived from MCQ; award credit for stating the correct idea and giving supporting reasoning.
\item \textbf{Answer:} Social functioning is a reflection of psychological functioning. Explain why this is the correct answer and provide supporting evidence. Reference key facts or concepts that support this choice.
\\ \textit{Note:} Long-form derived from MCQ; award credit for stating the correct idea and giving supporting reasoning.
\end{enumerate}

\vfill
\noindent\textit{End of Answer Key}
\end{document}